\chapter{An introduction to survival models}

\section{Survival function}


Survival analysis is also called time to event analysis or event
history analysis.  We use survival models to analyze data that have the following characteristics:

\begin{enumerate}
\item the dependent variable is the waiting time till the occurrence of some event;
\item some observations are censored; that is, some units are not in the data set anymore for some reason before an event happens.  For some units we observe the event and the exact waiting time; for others it has not occurred, and all we know is that the waiting time exceeds the observation time.  However, we believe if they stayed in the study, sooner or later an event would happen.
\item we are trying to explain (predict) the occurrence of an event with some predictors (explanatory variables).
\end{enumerate}

Let $T$ denote the time variable, a nonnegative, continuous
random variable with PDF $f(t)$ and CDF $F(t)$, where $t$ is a
realization of $T$.  The survival function is defined by

\begin{equation}
S(t) = Pr(T>t) = 1-F(t)
\end{equation}

The graph of $S(t)$ against $t$ is known as the survival curve.  If
there is no censoring, then the survival function is easy to estimate
by the ratio of number of units survives at $t$ vs. total number of
units at risk at $t$.  The most common method to estimate the survival
function with a survival data set containing censored observations is
the Kaplan-Meier method.  The basic idea is to use the product of a series of conditional probabilities.

First sort the event times from the smallest to the largest:


\begin{equation}
t_{(1)} \le t_{(2)} \le ... \le t_{(r)}
\end{equation}


Then the survival curve is estimated by:

\begin{equation}
\hat S(t) = \prod_{j | t{(j)} \le t} (1-\frac{d_j}{r_j}),
\end{equation}

where $r_j$ is the number of units at risk at $t_{(j)}$ and $d_j$ is
the number of ``failure'' or events at $t_{(t)}$.  Note that units
censored at $t_{(j)}$ are included in $r_j$.

The variance of this estimator can be estimated by:

\begin{equation}
\hat {var} (S(t)) = (\hat S(t))^2 \sum_{j | t{(j)} \le t} \frac{d_j}{r_j(r_j-d_j)}.
\end{equation}



\section{Hazard function}

We are often interested in which periods have the highest or lowest
change of ``death'' or ``failure'' or other events.  That is, we may
be interested in the probability that a state ends between $t$ and $t+
\Delta t$ conditional on having reached $t$ in the first place. 

The probability is
\begin{equation}
\mbox{Pr} (t<T \leq t+\Delta t | T \geq t)=\frac{F(t+\Delta
t)-F(t)}{S(t)}
\end{equation}

The hazard function is defined by
\begin{equation}
h(t)=\lim_{\Delta t \rightarrow 0}\mbox{Pr} (t<T \leq t+\Delta t |
T \geq t)=\frac{f(t)}{S(t)}=\frac{f(t)}{1-F(t)}
\end{equation}

One of the simplest functional forms is the exponential
distribution
\begin{equation}
f(t,\theta)=\theta e^{-\theta t}
\end{equation}

Therefore, the hazard function is
\begin{equation}
h(t)=\theta
\end{equation}

A much more flexible functional form is Weibull
\begin{equation}
f(t,\theta, \alpha)=1-\exp(-(\theta t)^{\alpha}).
\end{equation}

Hazard function can be shown to be
\begin{equation}
h(t)=\alpha \theta^{\alpha} t^{\alpha-1}
\end{equation}

However, Weibull distribution does not allow for the possibility
that the hazard may first increase then decrease over time.
Log-normal distribution does allow this.

% \begin{figure}[htbp] \footnotesize \caption{\small Various hazard functions from Davidson and MacKinnon}
% \centerline{\includegraphics[width=4in,height=3in]{davidson2.eps}}
% \label{fig:harzard}
% \end{figure}

\section{Log-likelihood function}

Suppose we have $n$ units with a survival function $S(t)$, density $f(t)$, and hazard $h(t)$.
If for unit $i$, the event happens at $t_i$, it's contribution to the likelihood function is the density at that point:

\begin{equation}
L_i=f(t_i)=S(t_i)h(t_i)
\end{equation}


If no event at $t_i$, all we know is that the lifetime exceeds $t_i$.  The probability is 
\begin{equation}
L_i=S(t_i), 
\end{equation}
which is the contribution to the likelihood function from a censored
observation. 

Therefore, if $U$ denotes the set of uncensored observations, the
log-likelihood function for the entire sample can be written as
\begin{equation}
\ell (t, \theta)=\sum_{i \in U} \log h(t_i | {\bf X}_i,
\theta)+\sum_{i=1}^n \log S(t_i | {\bf X}_i, \theta)
\end{equation}

\section{Proportional Hazard Models}

\subsection{Model and likelihood}

One class of models that is quite widely used is the proportional
hazard models (based on Cox (1972)).
\begin{equation}
h(t | {\bf x}_i)=h_0(t)\exp ({\bf x}'_i \beta),
\label{cox1}
\end{equation}

In this model $h_0(t)$ is called baseline hazard rate; it describes
the risk for units with $\bf x_i = 0$. $\exp ({\bf x}'_i \beta)$
represents the relative risk, a proportionate increase or decrease in
risk, due to $\bf x_i$.  The interpretation of $\beta$ is that
$\exp(\beta_i)$ gives the relative risk change associated with an
increase of one unit in $x_i$, all other explanatory variables remaining
constant.

$\beta$ is estimated by maximizing a partial likelihood function.  It
is partial likelihood function since the baseline hazard is factored
out.  Suppose one event happens at time $t_j$.  Conditional on this
event the probability that case $i$ dies (meaning this event happened
to case $i$, not other cases in the risk set) is


\begin{equation}
L_i=\frac{h_0(t) \exp(x_i' \beta)} {\sum_l I(T_l \ge t) h_0(t) \exp(x_l' \beta)}=\frac{ \exp(x_i' \beta)} {\sum_l I(T_l \ge t)  \exp(x_l' \beta)}
\end{equation}

We can see here that Cox's model is for continuous time survival data.
It is assumed there is no ties at time $t_j$ (in continuous time is
not possible for two events to happen at the same time).  In reality,
we observe ties of event time, because in many situations, we observe
grouped data.  For example, we observe case $m$ and $n$ have events at
the same time $j$.  In that case, in the denominator, whether we include
case $m$ in the calculation of the likelihood of case $n$ is
ambiguous.  In that case, special treatment is needed.  A standard
approximation (Breslow and Peto) is to let

\begin{equation}
L_{m \in D(t_j)} \simeq \frac{ \prod_{m \in D(t_j)} \exp(x_m' \beta)} {[\sum_{l \in R(t_j)} \exp(x_l' \beta)]^{d_j}}
\end{equation}

where $D(t_j)$ is the set of cases that die at time $t_j$ and $d_j$
denotes the number of cases that die at time $t_j$.  This
approximation works well with small number of ties relative to total
number of cases at risk.


One important feature (disadvantage?) for the proportional hazard
model is that for all units, the effect is the same at all time $t$.
To see this:

\begin{equation}
\frac{h_1(t)}{h_2(t)} = \frac{h_0(t)\exp ({\bf x_1}'_i \beta)}{h_0(t)\exp ({\bf x_2}'_i \beta)}= \frac{\exp ({\bf x_1}'_i \beta)}{\exp ({\bf x_2}'_i \beta)},
\end{equation}
which does not depend on $t$.  

One advantage of Cox's model is that it is considered as a
semi-parametric model since it does not have to specify the
distribution of the survival time.  The estimation process relies only
on the order in which events occur, not the exact time they occur.
Parameter estimates are derived assuming continuous survival times.

\subsection{Interpretation}

In equation \ref{cox1}, if we would like to calculate the partial effect of $x_j$ on $h$, we have:

\begin{equation}
\partial h(t | {\bf x}_i) / \partial x_j=h_0(t)\exp ({\bf x}'_i \beta) \beta_j = \beta_j h(t | {\bf x}_i),
\end{equation}

If we do it the other way, plug in $x_j + 1$ will generate

\begin{equation}
h_0(t)\exp ({\bf x}'_i \beta + \beta_j) = \exp (\beta_j) h(t|{\bf x}'_i \beta)
\label{cox2}
\end{equation}


Therefore, one unit change in $x_j$ will have a change of $1- \exp(\beta_j)$ times the original hazard.

\section{Discrete-time Survival Models}

Cox (1972)  proposed a discrete-time model by modeling the odds of dying at time $t_j$ given survival up to that point.

\begin{equation}
\frac{h(t_j | {\bf x}_i)}{1- h(t_j | {\bf x}_i)} = \frac{h_0(t_j | {\bf x}_i)}{1- h_0(t_j | {\bf x}_i)} \exp({\bf x}'_i \beta),
\end{equation}
which is similar to Cox's proportion hazard model for continuous time.  If we take logs on both sides, we have 

\begin{equation} 
{\rm logit} (h(t_j | {\bf x}_i)) = \alpha_j +  {\bf x}'_i \beta,
\end{equation}
This looks very similar to a logit model.  In fact, we can fit a
discrete-time proportional hazard model by running a logistic
regression on a set of pseudo observations generated by ``filling in''
the observations between the start time to the event time or time at
the end of the study period (censored).  The outcome of this process
is 0 unless there is an event, then it turns 1.  It is treated as
independent Bernoulli observations with probability given by hazard
$h_{ij}$ for individual $i$ at time $t_j$.  The likelihood function
for the discrete-time survival model coincide with the binomial
likelihood of independent binomial process (\cite{Sing:Will:it's:1993} has more details).

\section{Do it in Stata}

To do a discrete-time survival model in Stata, your data set needs to
be set up as by unit-time.  For example, observations by company
month.  The event (or failure) variable needs to be set to zero until
the event happens.  If no event until the end of the study period,
then it is censored.  A data set by unit-time can be analyzed by logit
model, optionally with clustered standard error.  To set up a
structure like this, user-written programs called $dthaz$ and
$prsnperd$ (means person-period) can be used.


To do a continuous-time survival model in Stata, your data can be
either set up as discrete-time case (in which case you can use a
time-varying explanatory variable) or a single-record-per-person data
set.  Stata has build-in command $stset$ to set up a survival analysis
structure.  Then you can run $stcox$ for Cox's proportional hazard
model.

% \bibliography{singer}
